%  USE CASE 1-1
\begin{figure}[H]
\centering
\begin{tikzpicture}[
    scale=0.9, transform shape,
    % --- STYLES ---
    action/.style={
        rectangle, rounded corners=0.3cm,
        draw=black, fill=white,
        minimum width=3.2cm, 
        minimum height=1cm,
        align=center, font=\small,
        inner sep=4pt
    },
    decision/.style={
        diamond, draw=black, fill=white,
        aspect=2, inner sep=2pt,
        font=\footnotesize, minimum width=2.5cm,
        align=center 
    },
    merge/.style={
        diamond, draw=black, fill=white,
        aspect=2, inner sep=2pt,
        minimum width=0.9cm,
        minimum height=0.9cm
    },
    start/.style={circle, fill=black, minimum size=0.5cm},
    end/.style={circle, draw=black, double, minimum size=0.5cm, inner sep=1pt},
    control/.style={->, thick, >=Stealth, rounded corners=5pt}
]

% ===== 1. COORDINATES & HEADERS =====
\def\xSales{0} 

\node[font=\bfseries\large] at (\xSales, 1) {Nhân viên Sales};

% ===== 2. SALES LANE NODES =====
\node[start] (start) at (\xSales, 0) {};

\node[action, below=0.8cm of start] (s1) {Gửi yêu cầu\\ thuê phòng};

\node[action, below=1cm of s1] (s2) {Tiếp nhận thông tin\\ nhu cầu khách hàng};

\node[decision, below=1cm of s2] (dec) {Phòng/giường có\\ đáp ứng đúng\\ nhu cầu khách hàng?};

\node[action, below=1cm of dec] (s3_yes) {Tư vấn các\\ loại dịch vụ};
\node[action, right=1.2cm of dec] (s3_no) {Gợi ý dịch vụ\\ khác của KTX};

\node[merge, below=1cm of s3_yes] (merge1) {};
\node[action, below=0.9cm of merge1] (s4) {Xác nhận nhu cầu\\ thuê phòng};
\node[action, below=1cm of s4] (s5) {Thông báo kết quả tiếp nhận\\ yêu cầu cho khách hàng};

\node[end, below=1cm of s5] (end) {};
\fill (end.center) circle (0.15cm);

% ===== 3. CONNECTIONS =====

% -- Control Flow --
\draw[control] (start) -- (s1);
\draw[control] (s1) -- (s2);
\draw[control] (s2) -- (dec);

% Decision splits
\draw[control] (dec) -- node[right, font=\footnotesize] {Có} (s3_yes);
\draw[control] (dec) -- node[above, font=\footnotesize] {Không} (s3_no);

% Merging back to step 4
\draw[control] (s3_yes) -- (merge1);
\draw[control] (s3_no.south) |- (merge1.east);

\draw[control] (merge1) -- (s4);
\draw[control] (s4) -- (s5);
\draw[control] (s5) -- (end);

\end{tikzpicture}
\caption{Activity Diagram UC1-1}
\end{figure}
